\documentclass{article}

% if you need to pass options to natbib, use, e.g.:
%     \PassOptionsToPackage{numbers, compress}{natbib}
% before loading neurips_2025

% The authors should use one of these tracks.
% Before accepting by the NeurIPS conference, select one of the options below.
% 0. "default" for submission
% \usepackage{neurips_2025}
% the "default" option is equal to the "main" option, which is used for the Main Track with double-blind reviewing.
% 1. "main" option is used for the Main Track
%  \usepackage[main]{neurips_2025}
% 2. "position" option is used for the Position Paper Track
%  \usepackage[position]{neurips_2025}
% 3. "dandb" option is used for the Datasets & Benchmarks Track
 % \usepackage[dandb]{neurips_2025}
% 4. "creativeai" option is used for the Creative AI Track
%  \usepackage[creativeai]{neurips_2025}
% 5. "sglblindworkshop" option is used for the Workshop with single-blind reviewing
 % \usepackage[sglblindworkshop]{neurips_2025}
% 6. "dblblindworkshop" option is used for the Workshop with double-blind reviewing
%  \usepackage[dblblindworkshop]{neurips_2025}

% After being accepted, the authors should add "final" behind the track to compile a camera-ready version.
% 1. Main Track
 % \usepackage[main, final]{neurips_2025}
% 2. Position Paper Track
%  \usepackage[position, final]{neurips_2025}
% 3. Datasets & Benchmarks Track
 % \usepackage[dandb, final]{neurips_2025}
% 4. Creative AI Track
%  \usepackage[creativeai, final]{neurips_2025}
% 5. Workshop with single-blind reviewing
%  \usepackage[sglblindworkshop, final]{neurips_2025}
% 6. Workshop with double-blind reviewing
%  \usepackage[dblblindworkshop, final]{neurips_2025}
% Note. For the workshop paper template, both \title{} and \workshoptitle{} are required, with the former indicating the paper title shown in the title and the latter indicating the workshop title displayed in the footnote.
% For workshops (5., 6.), the authors should add the name of the workshop, "\workshoptitle" command is used to set the workshop title.
% \workshoptitle{WORKSHOP TITLE}

% "preprint" option is used for arXiv or other preprint submissions
\usepackage[preprint]{neurips_2025}

% to avoid loading the natbib package, add option nonatbib:
%    \usepackage[nonatbib]{neurips_2025}

\usepackage[utf8]{inputenc} % allow utf-8 input
\usepackage[T1]{fontenc}    % use 8-bit T1 fonts
\usepackage{hyperref}       % hyperlinks
\usepackage{url}            % simple URL typesetting
\usepackage{booktabs}       % professional-quality tables
\usepackage{amsfonts}       % blackboard math symbols
\usepackage{nicefrac}       % compact symbols for 1/2, etc.
\usepackage{microtype}      % microtypography
\usepackage{xcolor}         % colors
\usepackage{caption}
\usepackage{subcaption}
\usepackage{algorithm}
\usepackage{algorithmic}
\usepackage{graphicx}
\usepackage{float}
\usepackage{amsmath}
% Note. For the workshop paper template, both \title{} and \workshoptitle{} are required, with the former indicating the paper title shown in the title and the latter indicating the workshop title displayed in the footnote. 
\title{NELE: Non-Uniform Rational B-Spline Elementwise Learnable}


% The \author macro works with any number of authors. There are two commands
% used to separate the names and addresses of multiple authors: \And and \AND.
%
% Using \And between authors leaves it to LaTeX to determine where to break the
% lines. Using \AND forces a line break at that point. So, if LaTeX puts 3 of 4
% authors names on the first line, and the last on the second line, try using
% \AND instead of \And before the third author name.


\author{%
  Vishakh.~Begari\thanks{Use footnote for providing further information
    about author (webpage, alternative address)---\emph{not} for acknowledging
    funding agencies.} \\
  Independent Researcher\\
  Bavaria, Germany \\
  \texttt{vibedjentle@gmail.com} \\
  % examples of more authors
  % \And
  % Coauthor \\
  % Affiliation \\
  % Address \\
  % \texttt{email} \\
  % \AND
  % Coauthor \\
  % Affiliation \\
  % Address \\
  % \texttt{email} \\
  % \And
  % Coauthor \\
  % Affiliation \\
  % Address \\
  % \texttt{email} \\
  % \And
  % Coauthor \\
  % Affiliation \\
  % Address \\
  % \texttt{email} \\
}


\begin{document}


\maketitle


\begin{abstract}
  The design of activation functions (AF) 
    is central to the expressiveness and efficiency of neural networks. 
    Yet the pursuit of a universal activation function, one that can 
    automatically adapt its shape to suit a wide variety of tasks, remains 
    an open challenge. In this work, the author proposes \texttt{NELE} (NURBS Elementwise Learnable), 
    a learnable activation function based on Non-Uniform Rational B-Splines (NURBS).
    Unlike conventional fixed activations, \texttt{NELE} 
    learns a smooth, data driven mapping from inputs to outputs 
    through a set of control points and corresponding weights.
    This yields a continuous, differentiable family of functions capable 
    of modeling a broad spectrum of nonlinear behaviors, effectively 
    unifying diverse activation patterns within a single formulation.
    \texttt{NELE} was evaluated on several state-of-the-art benchmark datasets.
    The empirical results indicate that \texttt{NELE} achieves consistently better performance 
    than multiple existing AF approaches. The theoretical universality of NURBS, its ability to approximate a wide class 
    of functions and its inherent flexibility, positions \texttt{NELE} as a promising candidate 
    for a universal activation function.

\end{abstract}
\section{Introduction}
Activation functions (AF) are a core component of artificial neural networks, determining how neurons process input 
signals and introduce non-linearity into the network. Their evolution reflects the growing complexity 
and depth of neural models: from the early binary step functions of McCulloch-Pitts \cite{mcculloch1943logical}
neurons and perceptrons, 
which could handle only simple, linearly separable tasks, to smooth, differentiable functions like \texttt{sigmoid} \cite{MINAI1993845}
and \texttt{tanh} \cite{6795724} that enabled backpropagation in multilayer networks. The advent of Rectified Linear Unit (\texttt{ReLU}) \cite{relu} and its variants solved 
vanishing gradient issues and facilitated deep learning, while modern activations like Exponential Linear Unit (\texttt{ELU}) \cite{clevert2016fastaccuratedeepnetwork}, 
\texttt{Swish} \cite{ramachandran2017searchingactivationfunctions}, 
\texttt{Mish} \cite{misra2020mishselfregularizednonmonotonic} 
further enhance network performance by improving training stability and modeling complex patterns.
While classical activations such as \texttt{ReLU}, \texttt{tanh}, and Gaussian Error Linear Unit (\texttt{GELU}) \cite{hendrycks2023gaussianerrorlinearunits}
have enabled remarkable success, each represents only a narrow subset of possible nonlinear behaviors. 
Polynomial activation functions are unbounded and prone to numerical instability, 
often requiring high degrees or large networks to achieve sufficient expressiveness \cite{Hornik1989Multilayer}
\cite{Leshno1993NonPolynomial} \cite{Goodfellow2016DeepLearning}.

The concept of spline-based AF is not new. Early studies, such as those by 
\cite{guarnieri1999multilayer}, introduced adaptive spline activations to provide flexible, 
data-driven nonlinearities within neural networks. 
More recent works, including \cite{hryniowski2018polyneuronautomaticneurondiscovery} and 
\cite{bohra2020learning}, extended this idea by learning spline parameters directly during training, 
allowing the AF itself to adapt to the data. 
Building on this foundation, Non-Uniform Rational B-Splines (NURBS) can be seen as a natural 
extension, offering even greater flexibility through rational weighting and non-uniform knot placement. 
This enables NURBS based activations to represent a broader class of nonlinear mappings while 
retaining the smoothness and local control properties of traditional splines.
Using appropriate parameters, existing activation functions can be approximated using the NURBS based approach.

\section{Theory}
NURBS \cite{nurbs} are a mathematical representation used to describe curves and surfaces 
in a flexible and precise way. They generalize B-splines and Bézier \cite{Prautzsch2002Bezier} curves 
by introducing non-uniform knot vectors and rational weighting, 
allowing the representation of both free-form shapes and 
exact analytic geometry. The non-uniform spacing of knots enables localized shape control, 
while rational weights allow exact representation of conic sections 
and other geometries not achievable with polynomial B-splines.
A NURBS curve of degree $p$ is defined as:
\begin{equation}\label{equation:nurbs}
\mathbf{C}(u)=
\frac{
\displaystyle\sum_{i=0}^{n} N_{i,p}(u)\, w_i\, \mathbf{P}_i
}{
\displaystyle\sum_{i=0}^{n} N_{i,p}(u)\, w_i
},
\qquad u \in [u_0, u_m].
\end{equation}
where, $P_i$ = control points, $w_i$ = weights associated with the control points, 
$N_{i,p}(u)$ = B-spline basis functions of degree $p$, ${u_0, u_1,...u_m}$ = knot vector, 
$n + 1$ = number of control points and $m = n + p + 1$ = knot count relationship.

Setting all the weights equal ($w_i$ = 1), Equation~\ref{equation:nurbs} reduces to 
$C(u) = \displaystyle\sum_{i=0}^{n} N_{i,p}(u)P_i$, which is a B-spline curve.
Stone-Weierstrass theorem \cite{stone1948generalized} establishes that any continuous real valued 
function on a compact interval can be uniformly approximated arbitrarily 
closely by elements of a suitably rich algebra of functions 
(for example, polynomials), spline spaces being piecewise polynomial 
subalgebras, inherit the same dense approximation property 
in the space of continuous functions. Since NURBS curves 
generalize B-splines by adding weights 
(and thereby include B-spline curves as a special case), the ability to represent exactly all 
polynomials up to a given degree and retain the capacity to approximate any 
continuous (and sufficiently smooth) function arbitrarily well.

For a smooth curve $y = f(x)$, the curvature is given by:
\begin{equation}
\kappa = \frac{\left|y''\right|}{\left(1 + (y')^2\right)^{3/2}}
\end{equation}
The maximum curvature occurs at points where $\frac{d\kappa}{dx} = 0$.
The solutions for $x$ corresponding to $\kappa' = 0$ for different AF are presented
in Table ~\ref{tab:AF}.
Since curvature is a geometric invariant
\cite{doCarmo1976}, let $A(u) = \displaystyle\sum_{i=0}^{n} N_{i,p}(u)\, w_i\, \mathbf{P}_i$, \quad 
$W(u) = \displaystyle\sum_{i=0}^{n} N_{i,p}(u)\, w_i$, then
\begin{equation}
\mathbf{C'}=\frac{AW' - A'W}{W^2}
\end{equation}
\begin{equation}
\mathbf{C''}=\frac{W(A''W - AW'') - 2W'(AW' - A'W)}{W^3}
\end{equation}
\begin{equation}\label{eq:nurbscurv}
  \kappa = \frac{W^3(W(A''W - AW'') - 2W'(AW' - A'W))}{(W^4 + (AW' - A'W)^2)^\frac{3}{2}}
\end{equation}
The upper bound for Equation~\ref{eq:nurbscurv} is obtained by setting the denominator to zero.
\begin{equation}
(W^4 + (AW' - A'W)^2)^\frac{3}{2} = 0
\end{equation}
\begin{equation}
  W^4 + (AW' - A'W)^2 = 0
\end{equation}
Sum of two non-negative numbers is zero only if each term is zero individually,
$W^4 = 0$, $(AW' - A'W)^2 = 0$. Simplifying gives, $W = 0$, and $A'W = AW'$. 
Substituting, $W = 0$, $AW' = 0$. In a generic cubic Bézier curve with arbitrary weights and distinct control points, 
$A \neq 0$, in general. In that case, for $n = 3$, $W'$ is linear. This implies that the upper bound of curvature, $\kappa$,
is directly proportional to the weight of the middle control point. Furthermore, if the 
x coordinates of the 3 control points are fixed, 
this results in a total of four unknown parameters that can be learned.
Figure ~\ref{fig:nurbs} illustrates how varying the weight of the middle control 
  point affects the curve: a lower weight ($w=0.2$) pulls the curve away 
  from the point, uniform weight ($w=1.0$) yields the standard quadratic B-spline, 
  and a higher weight ($w=5.0$) attracts the curve toward the control point. 
  Red markers indicate the control points used for all curves. 
\begin{figure}[H]
  \begin{center}
    \includegraphics[width=0.4\linewidth]{../../PICS/nurbs.pdf}
  \caption{NURBS curves generated from three control points with quadratic 
  degree ($p=2$) and different weights.}
  \label{fig:nurbs}
  \end{center}
\end{figure}
\begin{table}
    \centering
  \caption{AF with their peak curvature values and the input positions at which these maxima occur.}
  \label{tab:AF}
  \begin{tabular}{@{}lccc@{}}
    \toprule
    AF & $y$ & $x_{\kappa' = 0}$ & $\kappa_{\text{max}}$\\
    \midrule
    \texttt{Tanh}~\ref{app:tanh} & $\frac{e^{x} - e^{-x}}{e^{x} - e^{-x}}$ & 0.9196 & 0.5073\\
    \texttt{Sigmoid}~\ref{app:sigmoid} & $\frac{1}{1 + e^{-x}}$ & 1.3819 & 0.0924\\
    \texttt{Softplus}~\ref{app:softplus} & $\ln{({1 + e^{x}})}$ & -0.5 & 0.1913\\
    \texttt{ELU}~\ref{app:elu} & $
    \begin{cases}
      x & \text{if } x \geq 0, \\
      \alpha (e^x - 1) & \text{if } x < 0
    \end{cases}$ & $-\frac{1}{2} \ln(2 \alpha^2)$ & $\frac{2}{3\sqrt3}$\\
    \texttt{SiLU}~\ref{app:silu} & $x \cdot \frac{1}{1 + e^{-x}}$ & -0.3768 & 0.4037 \\
    \texttt{GELU}~\ref{app:gelu} & $x \phi(x)$ & -0.4773 & 1.3579\\
    \texttt{ReLU} & $
    \begin{cases}
    x, & \text{if } x > 0, \\
    0, & \text{if } x \le 0.
    \end{cases}$ & undefined & undefined\\
    \texttt{Leaky-ReLU} & $
    \begin{cases}
    x, & \text{if } x > 0, \\
    \alpha x, & \text{if } x \le 0.
    \end{cases}$ & undefined & undefined\\
    \texttt{LeLU}~\ref{app:LeLU} & $
    \begin{cases}
    x, & \text{if } x > 0, \\
    e^{(1 - \beta)x} - 1 + \beta x, & \text{if } x \le 0.
    \end{cases}$ & $- \frac{\sqrt{2}}{2}\ln2$ & 0.3924\\
    \texttt{Mish}~\ref{app:mish} & $x \cdot \tanh\!\left(\ln(1 + e^{x})\right)$ & -0.4271 & 0.5900\\
  \bottomrule
\end{tabular}
\end{table}
\section{Formulation}\label{sec:form}
The motivation for \texttt{NELE} was 
\begin{itemize}
\item Universality : To enable an AF which opens up a search 
space for flexibility.
\item Attenuation : Prevent large positive values from weakening \cite{elfwing2017sigmoidweightedlinearunitsneural}.
\item Non Zero-Gradient : Produce negative side activations.
\item Smoothnesss : Maintain continuity and differentiability throughout.
\end{itemize} 

The universality follows naturally from its 
NURBS-based approximation shown in Appendix~\ref{app:nurbsparams}. Attenuation prevention with non-zero gradients 
is achieved through a piecewise formulation in which the 
positive input region preserves the identity mapping of \texttt{ReLU}, 
while the negative region is shaped to emulate smooth activations such as 
\texttt{SiLU} or \texttt{GELU}. To ensure continuity and differentiability 
across the junction of these regions, the control points and their 
associated weights in the vicinity of the origin are appropriately adjusted, 
enforcing smooth transitions between the linear and nonlinear segments.

A cubic spline is employed as it provides sufficient expressivity while maintaining low 
computational complexity and numerical stability.
Four control points ($p_0$, $p_1$, $p_2$, $p_3$) are selected, corresponding to a cubic spline, 
which represents the lowest order capable of modeling nonlinear curvature while avoiding 
overfitting and their associated weights ($w_n$). The implementation is summarized in Algorithm ~\ref{alg:nele_training}.
Continuity condition forces $p_3 = [0.0, 0.0]$ and linearity of the positive side 
sets the slope between $p_2$ and $p_3$ = 1 to maintain tangency as depicted in Figure ~\ref{fig:nurbsregen}.
A mask is added which returns 0 for $x < p_0$.
The remaining control points and their associated weights were either treated as 
learnable parameters or fixed constants. 
For the fixed variant, setting $p_0$ = [-4.0, 0.0], $p_1 = [0.1, 0.1]$ and $w_n = 1$ yielded the best results. 
\begin{figure}[H]
  \begin{center}
    \includegraphics[width=0.4\linewidth]{../../PICS/Approx_gelu.pdf}
  \caption{GELU approximated for $x \leq 0$ using cubic NURBS with 4 control points and optimized weights.
  Range of NURBS curves used for hyperparameter tuning in experiments is also shown for comparison.}
  \label{fig:nurbsregen}
  \end{center}
\end{figure}
\begin{algorithm}
\caption{NELE activation function}
\label{alg:nele_training}
\begin{algorithmic}[1]
\REQUIRE Training data $\mathcal{D}$, initial parameters $p_0$, $p_1$, $p_2$, $p_3$, $w_0$, $w_1$, $w_2$, $w_3$
\ENSURE Optimized parameters
\IF{Learnable parameters}
    \STATE Replace $x_2$ and $y_2$ with $l/ \sqrt(2)$
    \STATE Initialize $l$, $p_0$, $p_1$, $p_3$
    \STATE Set $w_0$, $w_3$
\ELSE
    \STATE Set $p_0$, $p_1$, $p_2$, $p_3$, $w_0$, $w_1$, $w_2$, $w_3$
\ENDIF
\FOR{each batch}
    \IF{Clamped}
      \STATE Set $x_0$ = $min(\mathcal{D})$
    \ENDIF
    \STATE Evaluate $N_0 = (1-t)^3$, $N_1 = 3t(1-t)^2$, $N_2 = 3t^2(1-t)$, $N_3=t^3$, where $t \in [0,1]$
    \STATE Calculate $\text{Numerator\_y} = \sum N_n * w_n * y_n$, $\text{Numerator\_x} = \sum N_n * w_n * x_n$, $\text{Denominator} = \sum N_n * w_n$
    \STATE Calculate ${curve\_y} = \text{Numerator\_y}/\text{Denominator}$. ${curve\_x} = \text{Numerator\_x}/\text{Denominator}$.
    \STATE Find ${y}$ from ${curve\_y}$ for ${curve\_x}$ = $\mathcal{D}$ values via Linear Interpolation.
    \IF{Masking}
      \IF{Double}
        \IF{$min(\mathcal{D}) \leq\mathcal{D} \leq 0$}
        \RETURN $y$ 
        \ELSE
        \RETURN $\mathcal{D}$
        \ENDIF
        \ELSE
        \IF{$\mathcal{D} \leq 0$}
        \RETURN $y$ 
        \ELSE
        \RETURN $\mathcal{D}$
        \ENDIF
      \ENDIF
    \ELSE
      \RETURN $y$
    \ENDIF
\ENDFOR
\end{algorithmic}
\end{algorithm}
\section{Experiments with \texttt{NELE}}
Following the construction process outlined in Section ~\ref{sec:form}, 
\texttt{NELE} is benchmarked against a representative set 
of standard AF, \texttt{tanh}, \texttt{sigmoid}, \texttt{softplus}, 
\texttt{ELU}, \texttt{SiLU}, \texttt{GELU}, \texttt{ReLU}, 
\texttt{Leaky ReLU}, \texttt{LeLU} \cite{BIGARELLA2026122739} and 
\texttt{Mish}. Experiments are conducted on a combination of synthetic and standard benchmark datasets: 
(i)  Non-linear synthetic data (NLSD), (ii) MNIST \cite{lecun1998mnist}, (ii) MNIST autoencoder, (iv) CIFAR-10 \cite{hope2017learning}
 and (v) CIFAR-100 \cite{krizhevsky2009cifar}. 
All empirical evaluations are performed using feed-forward neural 
network architectures. The behavior of \texttt{NELE} in 
large-scale architectures, such as transformers or 
recurrent networks, is not investigated in this work.
All models are trained for a fixed number of epochs to ensure a 
fair comparison across AF under identical optimization conditions.
See Appendix~\ref{app:compute} for full details on all 
runs and total compute.
\subsection{Non-linear synthetic data (NLSD)}
NLSD provides a controlled environment for assessing functional expressivity.
The neural network used for curve fitting consists of a fully connected feed-forward 
architecture with a single hidden layer of 64 neurons. 
The hidden layer employs the various AF, while the output layer is linear. 
This design isolates the contribution of the activation function and 
ensures that observed differences in performance reflect the expressive 
capacity of each activation.
In schematic form, the network can be represented as:
\begin{displaymath}
x \;\xrightarrow{\text{Linear } (\text{input\_dim} \to 64)}\;
\text{AF} \;\xrightarrow{\text{Linear } (64 \to 1)}\; y
\end{displaymath}
The model was trained with Mean-Squared-Error (MSE) loss \cite{BickelDoksum2015MathStats}.
MSE provides a smooth, differentiable penalty for prediction errors, 
emphasizes larger deviations, and produces outputs that closely follow 
the target curve \cite{Bishop2006PRML}, which is important for subsequent smoothness analysis.
To evaluate the performance of different AF for function approximation, 
a series of synthetic datasets representing diverse nonlinear relationships is generated. 
Specifically, the target outputs $y$ were constructed as follows:
\begin{equation}
y_1 = \sin(x) + \epsilon\tag{a}
\end{equation}
\begin{equation}
y_2 = \sin(x) + 0.5 \cdot \sin(3x) + \epsilon\tag{b}
\end{equation}
\begin{equation}
y_3 = e^{-0.5 x} + \epsilon\tag{c}
\end{equation}
\begin{equation}
y_4 = \tanh(x) + \epsilon\tag{d}
\end{equation}
\begin{equation}
y_5 = x^2 +  \epsilon\tag{e}
\end{equation}
\begin{equation}
y_6 = \frac{x^3}{e^{x} - 1 + 10^{-6}} + \epsilon\tag{f}
\end{equation}
where $\epsilon \sim \mathcal{N}(0,1)$ denotes Gaussian noise added to simulate measurement 
or environmental variability. The chosen functions span a range of behaviors, 
including oscillatory patterns, exponential decay, saturating nonlinearities 
and polynomial growth, thereby providing a comprehensive testbed for assessing the 
generalization capability of different neural network architectures. 
The networks were evaluated based on their ability to fit the underlying 
function while mitigating overfitting to noise, $\epsilon = 0.2$, chosen relative to the scale of the input data. Each model was trained for 
a maximum of 300 epochs with a learning rate ($10^{-2}$, $10^{-3}$, $10^{-4}$) and parameter initialisation 
optimised to give the lowest training loss through grid search.
The studies were repeated across 7 runs with different random seeds to 
account for stochasticity in network initialization and data shuffling \cite{Dushatskiy_2022}.
and the median training loss across these runs is shown in 
Figure~\ref{fig:Lossplottrend}.
Preliminary results indicate that AF with smooth, bounded characteristics 
(e.g., \texttt{Tanh}, \texttt{Sigmoid}) excel in capturing oscillatory 
and saturating patterns.
These findings highlight the importance of matching AF properties 
to the underlying target function characteristics in curve-fitting tasks.
Given a set of discrete points $\{(x_i, y_i)\}_{i=1}^N$ representing a curve, 
smoothness was quantified using curvature based criteria:
For a parametric curve $(x(t), y(t))$, the curvature at a point is defined as
\begin{equation}
\kappa(t) = \frac{|x'(t)y''(t) - y'(t)x''(t)|}{\left( x'(t)^2 + y'(t)^2 \right)^{3/2}}
\end{equation}

The total curvature is approximated for discrete points using finite differences:
\begin{equation}
K_\text{total} = \sum_{i=1}^{N} \kappa_i
\end{equation}
\begin{figure}[H]
    \centering
    % Row 1
    \begin{subfigure}{0.32\textwidth}
        \centering
        \includegraphics[width=\linewidth]{../../PICS/NLSD/exp/curve_fitting.pdf}
        \caption{Exponential + Noise}
    \end{subfigure}
    \hfill
    \begin{subfigure}{0.32\textwidth}
        \centering
        \includegraphics[width=\linewidth]{../../PICS/NLSD/hyp/curve_fitting.pdf}
        \caption{Hyperbola + Noise}
    \end{subfigure}
    \hfill
    \begin{subfigure}{0.32\textwidth}
        \centering
        \includegraphics[width=\linewidth]{../../PICS/NLSD/quad/curve_fitting.pdf}
        \caption{Quadratic + Noise}
    \end{subfigure}

    % Row 2
    \begin{subfigure}{0.32\textwidth}
        \centering
        \includegraphics[width=\linewidth]{../../PICS/NLSD/sine/curve_fitting.pdf}
        \caption{Sinusoidal + Noise}
    \end{subfigure}
    \hfill
    \begin{subfigure}{0.32\textwidth}
        \centering
        \includegraphics[width=\linewidth]{../../PICS/NLSD/trig/curve_fitting.pdf}
        \caption{Trigonometric + Noise}
    \end{subfigure}
    \hfill
    \begin{subfigure}{0.32\textwidth}
        \centering
        \includegraphics[width=\linewidth]{../../PICS/NLSD/exppoly/curve_fitting.pdf}
        \caption{Polynomial + Noise}
    \end{subfigure}

    \caption{Curve fitting results for different target functions 
    using the proposed neural network, illustrating the network's 
    ability to approximate diverse nonlinear patterns.}
    \label{fig:curvefittingtrend}
\end{figure}

\begin{table}
\centering
  \caption{Comparison of smoothness metrics for different activation functions and \texttt{NELE} tuned for lowest training loss. 
The table reports total curvature ($K_\text{total}$) for each AF. 
Lower values indicate smoother behavior.}
  \label{tab:smooth}
  \begin{tabular}{@{}llllllll@{}}
    \toprule
    AF & (a) & (b) & (c) & (d) & (e) & (f) \\
    \midrule
    \texttt{Tanh} & 46.72 & 39.51 & 41.47 & 73.0 & 155.65 & 35.13 \\ 
    \texttt{Sigmoid} & 26.86 & 31.5 & 44.09 & 63.14 & 68.48 & 27.32 \\ 
    \texttt{Softplus} & 31.96 & 34.94 & 39.45 & 55.88 & 57.04 & 23.04 \\ 
    \texttt{ELU} & 45.33 & 33.36 & 39.77 & 57.36 & 61.63 & 23.44 \\ 
    \texttt{SiLU} & 46.02 & 32.34 & 39.51 & 57.38 & 64.14 & 30.44 \\ 
    \texttt{GELU} & 77.24 & 34.76 & 41.84 & 61.93 & 107.0 & 50.47 \\ 
    \texttt{ReLU} & 107.69 & 55.85 & 51.58 & 75.75 & 79.89 & 51.41 \\ 
    \texttt{Leaky-ReLU} & 96.3 & 52.63 & 47.37 & 66.68 & 79.93 & 34.55 \\ 
    \texttt{LeLU} & 27.58 & 31.24 & 39.82 & 69.49 & 111.18 & 19.86 \\ 
    \texttt{Mish} & 60.13 & 34.1 & 43.28 & 59.05 & 60.72 & 28.71 \\ 
    \texttt{NELE} & 44.02 & 38.33 & 41.4 & 61.62 & 114.7 & 30.81 \\ 
  \bottomrule
\end{tabular}
\end{table}

Among the AF evaluated, the proposed \texttt{NELE} activation 
consistently demonstrates comparable and sometimes superior performance across all benchmark functions. 
Quantitatively, networks employing \texttt{NELE} achieve lower training loss and smaller 
standard deviation of curve-fitting errors relative to alternative activations, 
indicating both improved accuracy and enhanced reliability.

Following the recommended procedure for \texttt{LELU}, the neural network should 
be trained via backpropagation to minimize the mean absolute error (MAE) 
between the predicted values and the true function at the training points, 
followed by computation of diffusion terms and application of the diffusion loss. 
In the experiments, however, only the standard MSE loss was used. 
Under these conditions, \texttt{NELE} produced  
predictions with smoothness comparable to \texttt{LeLU} and training 
loss values that were lower or similar across all benchmark functions.
\begin{figure}[H]
    \centering

    % Row 1
    \begin{subfigure}{0.32\textwidth}
        \centering
        \includegraphics[width=\linewidth]{../../PICS/NLSD/exp/Loss_Bar_Chart_with_Error_Bars.pdf}
        \caption{Exponential + Noise}
    \end{subfigure}
    \hfill
    \begin{subfigure}{0.32\textwidth}
        \centering
        \includegraphics[width=\linewidth]{../../PICS/NLSD/hyp/Loss_Bar_Chart_with_Error_Bars.pdf}
        \caption{Hyperbola + Noise}
    \end{subfigure}
    \hfill
    \begin{subfigure}{0.32\textwidth}
        \centering
        \includegraphics[width=\linewidth]{../../PICS/NLSD/quad/Loss_Bar_Chart_with_Error_Bars.pdf}
        \caption{Quadratic + Noise}
    \end{subfigure}

    % Row 2
    \begin{subfigure}{0.32\textwidth}
        \centering
        \includegraphics[width=\linewidth]{../../PICS/NLSD/sine/Loss_Bar_Chart_with_Error_Bars.pdf}
        \caption{Sinusoidal + Noise}
    \end{subfigure}
    \hfill
    \begin{subfigure}{0.32\textwidth}
        \centering
        \includegraphics[width=\linewidth]{../../PICS/NLSD/trig/Loss_Bar_Chart_with_Error_Bars.pdf}
        \caption{Trigonometric + Noise}
    \end{subfigure}
    \hfill
    \begin{subfigure}{0.32\textwidth}
        \centering
        \includegraphics[width=\linewidth]{../../PICS/NLSD/exppoly/Loss_Bar_Chart_with_Error_Bars.pdf}
        \caption{Polynomial + Noise}
    \end{subfigure}

    \caption{Standard deviation of prediction errors for different target functions, illustrating 
    the variability in the neural network's curve fitting performance. Despite \texttt{LeLU} 
    reaching a lower training loss for the sinusoidal function, 
    the corresponding fit lacks smoothness.}
    \label{fig:standarddevtrend}
\end{figure}
\begin{figure}[H]
    \centering

    % Row 1
    \begin{subfigure}{0.32\textwidth}
        \centering
        \includegraphics[width=\linewidth]{../../PICS/NLSD/exp/training_loss.pdf}
        \caption{Exponential + Noise}
    \end{subfigure}
    \hfill
    \begin{subfigure}{0.32\textwidth}
        \centering
        \includegraphics[width=\linewidth]{../../PICS/NLSD/hyp/training_loss.pdf}
        \caption{Hyperbola + Noise}
    \end{subfigure}
    \hfill
    \begin{subfigure}{0.32\textwidth}
        \centering
        \includegraphics[width=\linewidth]{../../PICS/NLSD/quad/training_loss.pdf}
        \caption{Quadratic + Noise}
    \end{subfigure}

    % Row 2
    \begin{subfigure}{0.32\textwidth}
        \centering
        \includegraphics[width=\linewidth]{../../PICS/NLSD/sine/training_loss.pdf}
        \caption{Sinusoidal + Noise}
    \end{subfigure}
    \hfill
    \begin{subfigure}{0.32\textwidth}
        \centering
        \includegraphics[width=\linewidth]{../../PICS/NLSD/trig/training_loss.pdf}
        \caption{Trigonometric + Noise}
    \end{subfigure}
    \hfill
    \begin{subfigure}{0.32\textwidth}
        \centering
        \includegraphics[width=\linewidth]{../../PICS/NLSD/exppoly/training_loss.pdf}
        \caption{Polynomial + Noise}
    \end{subfigure}

    \caption{Training loss curves for the neural network across different target functions, 
    showing the convergence behavior and learning dynamics for each 
    functional form and AF.}
    \label{fig:Lossplottrend}
\end{figure}

\subsection{MNIST}
The performance of various activation functions was evaluated by 
replicating the experimental setup of \cite{hendrycks2023gaussianerrorlinearunits}.
Fully connected neural networks with eight hidden layers of 128 neurons each were trained
using AF listed in Table~\ref{tab:smooth}. 
All networks were trained using the Adam optimizer with cross-entropy loss 
\cite{mao2023crossentropylossfunctionstheoretical} and
learning rates of $10^{-3}$, $10^{-4}$, and $10^{-5}$. Validation used $5,000$ samples 
to select the best model. In total, seven runs were performed for 75 epochs on the complete 
dataset with the optimized hyperparameters and their median values are used for comparison.
\texttt{NELE} shows comparable convergence at epoch 50.
Figure~\ref{fig:mnistloss} shows how each AF influences convergence speed and final loss.
\begin{figure}[H]
  \begin{center}
    \includegraphics[width=1.0\linewidth]{../../PICS/MNIST/training_loss.pdf}
  \caption{Training and test accuracy trends for MNIST over epochs using different AF.
  Shaded areas indicate ±1 standard deviation across runs for \texttt{NELE}\label{fig:mnistloss}}
  \end{center}
\end{figure}
\begin{figure}[H]
  \begin{center}
    \includegraphics[width=0.5\linewidth]{../../PICS/MNIST/noise_study.pdf}
  \caption{Robustness results for MNIST over noise strengths using different AF.\label{fig:mnistnoise}}
  \end{center}
\end{figure}
\texttt{NELE} converged within 30 epochs with a 
learning rate (\textit{lr}) of $0.001$ and achieved the lowest test error~\ref{tab:mnisttesult}. 
The trained model was also tested for robustness against
inputs containing uniform noise Unif[-a,a] added to each example at different levels a, where
a = 3 is the greatest noise strength. 
Here, \texttt{NELE} shows good robustness with its performace matching \texttt{GELU} and exceeding
other AF.
\begin{table}
    \centering
  \caption{Comparison of training loss and test error on the MNIST classification task
trained with different AF. \texttt{NELE} achieves the 
lowest test error among all activations evaluated.}
  \label{tab:mnisttesult}
  \begin{tabular}{@{}llcc@{}}
    \toprule
    AF &Train & Test (Noise=0) & Test (Noise=3)\\
    \midrule
    \texttt{Tanh} & 0.0076 & 2.7 & 22.94\\
    \texttt{Sigmoid} & 0.0072 & 4.45 & 16.47\\
    \texttt{Softplus}& 0.004 & 2.52 & 22.37\\
    \texttt{ELU}& 0.0056 & 2.44 & 21.51\\
    \texttt{SiLU} & 0.0045 & 2.28 & 24.64\\
    \texttt{GELU} & 0.0049 & 2.32 & 24.89\\
    \texttt{ReLU} & 0.0054 & 2.32 & 23.02\\
    \texttt{Leaky-ReLU}& 0.0044 & 2.35 & 21.22\\
    \texttt{LeLU} & 0.0002 & 2.12 & 18.5\\
    \texttt{Mish} & 0.0056 & 2.32 & 24.31\\
    \texttt{NELE}& 0.0056 & \textbf{1.83} & \textbf{25.49}\\
  \bottomrule
\end{tabular}
\end{table}

\subsection{MNIST autoencoder}
The neural network is a deep autoencoder from \cite{desjardins2015naturalneuralnetworks}
with Adam optimizer and a batch size of 64. A learning rate of $10^{-3}$ with Mean Squared Error loss 
produced the highest validation score with 5,000 samples. Figure~\ref{fig:mnistenc} shows that
with settings of \texttt{NELE} from Figure~\ref{fig:nurbsregen}, the AF
can perform quite well in tests.
\begin{figure}[H]
  \begin{center}
    \includegraphics[width=1.0\linewidth]{../../PICS/MNIST_ENC/training_loss.pdf}
  \caption{Training loss and test reconstruction error 
  trends for MNIST autoencoder over epochs using 
  different AF. Shaded areas indicate ±1 standard deviation across runs for \texttt{NELE}\label{fig:mnistenc}}
  \end{center}
\end{figure}

\subsection{CIFAR-10}
The neural network is a 9-layer network with the architecture and training
procedure from \cite{hendrycks2023gaussianerrorlinearunits}. The tuning was done 
using learning rates of $10^{-3}$, $10^{-4}$, and $10^{-5}$ for 100 epochs.
The best model was selected based on validation accuracy using 90-10 training/validaition split.
Seven runs were performed and the median values of 3 runs are reported.
Figure~\ref{fig:cifar10loss} shows the training loss and test accuracy calculated every 5 epochs for CIFAR-10 over epochs using different AF.

\begin{figure}[H]
  \begin{center}
    \includegraphics[width=1.0\linewidth]{../../PICS/CIFAR10/training_loss.pdf}
  \caption{Training and test accuracy trends for CIFAR10 over epochs using different AF.}
  \label{fig:cifar10loss}
  \end{center}
\end{figure}

\subsection{CIFAR-100}
The model for CIFAR-100 is a wide residual neural network described in 
\cite{zagoruyko2017wideresidualnetworks} with learning rate scheduler and dropout keep
probability from
\cite{loshchilov2017sgdrstochasticgradientdescent} and batch normalization 
at the end of a residual block to stabilize exploding ELU's \cite{Shah_2016}
as implemented in \cite{hendrycks2023gaussianerrorlinearunits}. Setting from Figure~\ref{fig:nurbsregen} 
were used to obtain median convergence
curves over three runs as shown in Figure~\ref{fig:cifar100loss}.
\begin{figure}[H]
  \begin{center}
    \includegraphics[width=1.0\linewidth]{../../PICS/CIFAR100/training_loss.pdf}
  \caption{Training and test accuracy trends for CIFAR100 over 50 epochs using different AF.Shaded areas indicate ±1 standard deviation across runs for \texttt{NELE}}
  \label{fig:cifar100loss}
  \end{center}
\end{figure}
\section{Limitations}

\section{Conclusions}
\texttt{NELE} exhibits a combination of properties that make it particularly 
well suited for functional approximation tasks. 
Its inherent smoothness ensures that gradients are well behaved, 
facilitating stable optimization during training, 
while local control allows for fine grained adjustments in specific 
regions of the input space without affecting the entire function. 
In addition, geometric interpretability provides insight into 
how the function behaves across different regions and its capacity to 
represent both polynomial and non-polynomial shapes with a relatively 
small number of control points makes it highly efficient. 
Due to its theoretical universality, \texttt{NELE} can, in principle, 
approximate any continuous function to arbitrary accuracy, 
making it a flexible candidate for general-purpose activation in neural networks. 
To evaluate its practical performance, \texttt{NELE} was tested across multiple 
state of the art benchmark datasets spanning diverse machine 
learning tasks. The empirical results consistently demonstrate that 
\texttt{NELE} achieves superior performance relative to existing activation 
function methods, both in terms of convergence speed and predictive accuracy, 
highlighting its potential as a universal activation function capable of 
capturing complex functional relationships.

\begin{ack}
This work was carried out independently by the author.
\end{ack}

\bibliographystyle{unsrt}
\bibliography{NELE}


%%%%%%%%%%%%%%%%%%%%%%%%%%%%%%%%%%%%%%%%%%%%%%%%%%%%%%%%%%%%

\appendix
\section{Function analysis}
\begin{figure}[H]\label{fig:afcollect}
  \begin{center}
    \includegraphics[width=0.7\linewidth]{../../analysis/analysis.pdf}
  \caption{The solid lines represent the functions $y(x)$ for different activation functions. 
The points of maximum curvature, $\kappa_{\mathrm{max}}$, are indicated by markers along each curve. 
The dashed lines represent the curvature, $\kappa(x)$, as a function of $x$. 
Each color corresponds to a different activation function as shown in the legend.}
  \end{center}
\end{figure}
\subsection{Tanh}\label{app:tanh}

\begin{displaymath}
  y = \tanh{x}
\end{displaymath}
\begin{displaymath}
y' = \operatorname{sech}^2(x) = 1 - y^2
\end{displaymath}
\begin{displaymath}
y'' = -2\operatorname{sech}^2(x)\tanh{x} = -2y(1 - y^2)
\end{displaymath}
\begin{displaymath}
\kappa = \frac{|-2\operatorname{sech}^2(x)\tanh{x}|}{(1 + (\operatorname{sech}^2(x))^2)^\frac{3}{2}}
\end{displaymath}
\begin{displaymath}
\frac{d \kappa}{dx} = 0 
\end{displaymath}

\begin{displaymath}
\text{Using chain rule}, \frac{d \kappa}{dx} = \frac{d \kappa}{dy}\frac{dy}{dx}\text{. Then } \frac{d \kappa}{dy} = 0 \text{ or } 1 - y^2 = 0
\end{displaymath}
\begin{displaymath}
  \frac{d}{dy}\left( \frac{2y(1 - y^2)}{(1 + (1 - y^2)^2)^{3/2}}\right) = 3y^6 - 5y^4 - 2y^2 + 2 = 0
\end{displaymath}
\begin{displaymath}
  \text{Solving: }y = 0.9196, \kappa_{\max} = 0.5073
\end{displaymath}

\subsection{Sigmoid}\label{app:sigmoid}

\begin{displaymath}
  y = \frac{1}{1 + e^{-x}}
\end{displaymath}
\begin{displaymath}
  y' = \frac{-e^{-x}}{(1 + e^{-x})^2}
\end{displaymath}
\begin{displaymath}
  y'' = \frac{e^{-x}(1 - e^{-x})}{(1 + e^{-x})^3}
\end{displaymath}
\begin{displaymath}
  \kappa = \frac{\frac{e^{-x}(1 - e^{-x})}{(1 + e^{-x})^3}}{(1 + ( \frac{-e^{-x}}{(1 + e^{-x})^2})^2)^\frac{3}{2}}
\end{displaymath}
Solving this equation leads to a transcendental condition, 
which has no simple closed-form solution, so the 
locations of maximum curvature must be obtained numerically.

\subsection{Softplus}\label{app:softplus}

\begin{displaymath}
  y = \ln{({1 + e^{x}})}
\end{displaymath}
\begin{displaymath}
  y' = \frac{e^x}{1 + e^x}
\end{displaymath}
\begin{displaymath}
  y'' = \frac{e^{x}}{(1 + e^{x})^2}
\end{displaymath}
\begin{displaymath}
  \kappa = \frac{\frac{e^{x}}{(1 + e^{x})^2}}{(1 + (\frac{e^x}{1 + e^x})^2)^\frac{3}{2}}
\end{displaymath}
\begin{displaymath}
  \text{Let }s = \frac{e^x}{1 + e^x}, \kappa = \frac{s(1 - s)}{(1 + s^2)^\frac{3}{2}}
\end{displaymath}
\begin{displaymath}
  \frac{d\kappa}{dx} = 0, \frac{d}{dx}\bigg(\frac{s(1 - s)}{(1 + s^2)^\frac{3}{2}}\bigg) = 0
\end{displaymath}
\begin{displaymath}
  s = -1, 0.38197, 2.6180 \text{. Substituting, }x = -0.5, \kappa_{\text{max}} = 0.1913
\end{displaymath}

\subsection{ELU}\label{app:elu}
 \begin{displaymath}
  y = \begin{cases}
  x & \text{if } x \geq 0, \\
  \alpha (e^x - 1) & \text{if } x < 0
\end{cases}
 \end{displaymath}

 \begin{displaymath}
  y' = \begin{cases}
  1 & \text{if } x \geq 0, \\
  \alpha e^x & \text{if } x < 0
\end{cases}
 \end{displaymath}
\begin{displaymath}
  y'' = \begin{cases}
  0 & \text{if } x \geq 0, \\
  \alpha e^x & \text{if } x < 0
\end{cases}
\end{displaymath}
\begin{displaymath}
  \kappa = \begin{cases}
  0 & \text{if } x \geq 0, \\
  \frac{\alpha e^x}{(1 + \alpha^2 e^{2x})^\frac{3}{2}}  & \text{if } x < 0
\end{cases}
\end{displaymath}
\begin{displaymath}
  \frac{d \kappa}{dx} = 0
\end{displaymath}
\begin{displaymath}
  \frac{d}{dx} \left( \frac{\alpha e^x}{(1 + \alpha^2 e^{2x})^{3/2}} \right)
= \frac{\alpha e^x \left( 1 - 2 \alpha^2 e^{2x} \right)}{(1 + \alpha^2 e^{2x})^{5/2}} = 0
\end{displaymath}
\begin{displaymath}
  \alpha e^x \left( 1 - 2 \alpha^2 e^{2x} \right) = 0
\end{displaymath}
\begin{displaymath}
  e^x = 0 \text{ has no solution} \text{, therefore, } x = -\frac{1}{2} \ln(2 \alpha^2)
\end{displaymath}
\begin{displaymath}
  \text{Substituting back, }\kappa_{\text{max}} = \frac{2}{3\sqrt3}
\end{displaymath}

\subsection{SiLU}\label{app:silu}

\begin{displaymath}
y = \frac{x}{1 + e^{-x}}
\end{displaymath}
\begin{displaymath}
y' = \frac{1 + e^{-x} + xe^{-x}}{(1 + e^{-x})^2}
\end{displaymath}
\begin{displaymath}
y'' = \frac{e^{-x}(2 - x) + e^{-2x}(x+2)}{(1 + e^{-x})^3}
\end{displaymath}
\begin{displaymath}
\kappa = \frac{\frac{e^{-x}(2 - x) + e^{-2x}(x+2)}{(1 + e^{-x})^3}}{(1 + (\frac{1 + e^{-x} + xe^{-x}}{(1 + e^{-x})^2})^2)^\frac{3}{2}}
= \frac{(1 + e^{-x})^2(e^{-x}(2 - x) + e^{-2x}(x+2))}{[(1 + e^{-x})^4 + (1 + e^{-x} + xe^{-x})^2]^\frac{3}{2}}
\end{displaymath}
\begin{displaymath}
\frac{d \kappa}{dx} = 0
\end{displaymath}
\begin{displaymath}
\frac{d}{dx}\biggr(\frac{(1 + e^{-x})^2(e^{-x}(2 - x) + e^{-2x}(x+2))}{[(1 + e^{-x})^4 + (1 + e^{-x} + xe^{-x})^2]^\frac{3}{2}}\biggr) = 0
\end{displaymath}
Using quotient rule and consider the first term in the numerator after differentiation :
$[e^{-x}(x-3)-e^{-2x}(2x+3)](1+e^{-x})^2$. This leads to a nonlinear 
transcendental equation. Since, they do not simplify into algebraic 
forms that allow symbolic isolation of x, an exact analytical solution for the location of 
maximum curvature is not possible and numerical or graphical methods must be used.

\subsection{GELU}\label{app:gelu}

\begin{displaymath}
y = x \phi(x) \approx 0.5x(1 + \tanh [\sqrt{2 / \pi}(x + 0.44715 x^3)])
\end{displaymath}
\begin{displaymath}
\text{Let }k = \sqrt{\frac{2}{\pi}}, a = 0.441715 \text{ and }u(x) = k(x + ax^3)
\end{displaymath}
\begin{displaymath}
y' = 0.5(1 + \tanh{u}) + 0.5ax\operatorname{sech}^2(u)(1 + 3ax^2)
\end{displaymath}
\begin{displaymath}
y'' = 0.5a(1 + 3ax^2)\operatorname{sech}^2(u) + 0.5x[-2\operatorname{sech}^2(u)\tanh{u}(a(1 + 3ax^2))^2 + 6kax\operatorname{sech}^2(u)]
\end{displaymath}

\subsection{Mish}\label{app:mish}
\begin{displaymath}
y = x \cdot \tanh\!\left(\ln(1 + e^{x})\right)
\end{displaymath}
\begin{displaymath}
y' = \tanh{(\ln{(1 + e^x)})} + \frac{xe^x}{1 + e^x}\operatorname{sech}^2(\ln{(1 + e^x)})
\end{displaymath}
\begin{displaymath}
  y'' = a + b + c + d
\end{displaymath}
\begin{displaymath}
  \text{where, } a = \frac{e^x}{1 + e^x}\operatorname{sech}^2(\ln{(1 + e^x)})
\end{displaymath}
\begin{displaymath}
  b = \frac{\operatorname{sech}^2(\ln{(1 + e^x)})}{1+e^{-x}} 
\end{displaymath}
\begin{displaymath}
  c = \frac{xe^{-x}}{(1 + e^{-x})^2} \operatorname{sech}^2(\ln{(1 + e^x)})
\end{displaymath}
\begin{displaymath}
  d = \frac{-2x}{(1 + e^{-x})^2} \operatorname{sech}^2(\ln{(1 + e^x)} \tanh{\ln{(1 + e^x)}}
\end{displaymath}

\subsection{LeLU}\label{app:LeLU}

\begin{displaymath}
y = e^{(1 - \beta)x} + 1 - \beta x
\end{displaymath}
\begin{displaymath}
y' = (1 - \beta)e^{(1 - \beta)x} - \beta
\end{displaymath}
\begin{displaymath}
y'' = (1 - \beta)^2e^{(1 - \beta)x}
\text{Let }1 - \beta = a, e^{(1 - \beta)x} = e^{ax} = b
\end{displaymath}
\begin{displaymath}
\kappa = \frac{a^2b}{(1 + a^2b^{2})^\frac{3}{2}}
\end{displaymath}
\begin{displaymath}
\frac{d \kappa}{dx} = 0
\end{displaymath}
\begin{displaymath}
\text{Solving, }\beta = 1 \pm \sqrt{\frac{1}{2}}, 1, \text{Substituting }x = 0, \pm \frac{\sqrt{2}}{2}\ln2
\end{displaymath}

\section{Curve approximation with \texttt{NELE}\label{app:nurbsparams}}
To approximate a target curve with \texttt{NELE}, an optimization problem is formulated to minimize the mean squared error (MSE) between the target function and the NURBS representation.
Optimization algorithm L-BFGS-B \cite{scipy_optimize_minimize} is used to efficiently navigate the parameter space, 
producing a smooth curve that closely follows the desired shape.
\begin{table}[H]
  \caption{Control point coordinates and their respective weights for approximation on the negative space of AF using \texttt{NELE}}
  \label{tab:curveapprox}
  \centering
\scriptsize
  \begin{tabular}{@{}llllll@{}}
    \toprule
    AF&$p_0$&$p_1$&$p_2$&$p_3$&$w$\\
    \midrule
    \texttt{Tanh} & [-4.0, -1.0] &    [-0.884, -0.983] &    [-1.099, -1.078] &    [0.0, -0.001] &    [1.942 , 0.831 , 0.253 , 0.209]\\
    \texttt{Sigmoid} & [-4.0, 0.017] &    [-2.03, 0.059] &    [-1.498, 0.112] &    [0.0, 0.501] &    [1.31 , 0.886 , 0.719 , 0.808]\\
    \texttt{Softplus} & [-4.0, 0.017] &    [-2.332, 0.056] &    [-1.157, 0.105] &    [0.0, 0.694] &    [1.233 , 0.985 , 0.808 , 0.701]\\
    \texttt{ELU} & [-4.0, -0.984] &    [-2.394, -0.935] &    [-0.913, -0.912] &    [0.001, -0.001] &    [1.227 , 1.017 , 0.79 , 0.54]\\
    \texttt{SiLU} & [-4.0, -0.078] &    [-2.383, -0.112] &    [-1.011, -0.547] &    [0.0, 0.003] &    [1.226 , 1.012 , 0.806 , 0.611]\\
    \texttt{LeLU} & [-4.0, -1.373] &    [-2.375, -1.166] &    [-0.898, -0.917] &    [-0.001, 0.002] &    [1.235 , 1.015 , 0.776 , 0.524]\\
    \texttt{Mish} & [-4.0, -0.079] &    [-2.358, -0.113] &    [-0.868, -0.598] &    [-0.001, 0.008] &    [1.275 , 1.037 , 0.779 , 0.509]\\
    \texttt{GELU} & [-4.0, 0.0] &    [-0.321, 0.058] &    [-2.134, -0.913] &    [0.0, 0.0] &    [1.980 , 0.717 , 0.124 , 0.198]\\
  \bottomrule
\end{tabular}
\end{table}

\section{Online Resources}

The source code used in this paper are publicly available to ensure reproducibility and facilitate further research:

\begin{itemize}
    \item \textbf{\texttt{NELE} Implementation:} Python implementation of the proposed NELE activation function, including training scripts and examples. Available at: 
    \url{https://github.com/DjentleViBe/NELE}.
    
    \item \textbf{Additional Figures:} Addiitonal plots for the curve approximation from Table ~\ref{tab:curveapprox} are provided in the repository.
    \item \textbf{Hyperparameter tuning: }Parameters used in the analysis of NLSD were optimised through grid search with points 
    distributed equally from an initialised \texttt{GELU} approximation. Since grid search was unfeasible for other studies due to total run times, 
    hyperparameters were optimized using an automated search procedure. 
    For each candidate configuration, training was performed for a fixed number of 
    epochs with patience and validation performance evaluated after each epoch. Poorly performing 
    configurations were terminated early based on intermediate validation error, 
    and the objective value was defined as the minimum validation error attained during training after 200 trials using \cite{optuna_2019}.
    \item \textbf{Settings:} Parameters used for reproducing the results in this paper is provided in the
    \texttt{README} file of the repository.
  
\end{itemize}

It is recommended that users follow the provided \texttt{README} instructions for replicating experiments. This repository will remain publicly accessible for the foreseeable future.

\section{Compute Resources for Experiments}
\label{app:compute}
Experiments were conducted on an NVIDIA GTX 1080 Ti GPU 
with 64GB RAM.
\begin{table}[h]
\small
\centering
\begin{tabular}{l l l l l}
\textbf{Dataset} & \textbf{Time per iteration} & \textbf{Time per run} & \textbf{Runs} & \textbf{Total compute} \\
\hline
NLSD & ~0.003 s & ~0.2 s & 7 & ~2 s \\
MNIST & ~5 s & ~6.25 min & 7 & ~43 min \\
MNIST\_ENC & ~17 s & ~21 min & 7 & ~2.5 hr \\
CIFAR-10 & ~49 s & ~61 min & 3 & ~3 hr \\
CIFAR-100 & ~108 s & ~90 min & 3 & ~5 hr \\
\end{tabular}
\caption{Compute resources, per-run times, number of runs, and total compute for all experiments.}
\end{table}
%%%%%%%%%%%%%%%%%%%%%%%%%%%%%%%%%%%%%%%%%%%%%%%%%%%%%%%%%%%%

\newpage
\section*{NeurIPS Paper Checklist}

% %%% BEGIN INSTRUCTIONS %%%
% The checklist is designed to encourage best practices for responsible machine learning research, addressing issues of reproducibility, transparency, research ethics, and societal impact. Do not remove the checklist: {\bf The papers not including the checklist will be desk rejected.} The checklist should follow the references and follow the (optional) supplemental material.  The checklist does NOT count towards the page
% limit. 

% Please read the checklist guidelines carefully for information on how to answer these questions. For each question in the checklist:
% \begin{itemize}
%     \item You should answer \answerYes{}, \answerNo{}, or \answerNA{}.
%     \item \answerNA{} means either that the question is Not Applicable for that particular paper or the relevant information is Not Available.
%     \item Please provide a short (1–2 sentence) justification right after your answer (even for NA). 
%    % \item {\bf The papers not including the checklist will be desk rejected.}
% \end{itemize}

% {\bf The checklist answers are an integral part of your paper submission.} They are visible to the reviewers, area chairs, senior area chairs, and ethics reviewers. You will be asked to also include it (after eventual revisions) with the final version of your paper, and its final version will be published with the paper.

% The reviewers of your paper will be asked to use the checklist as one of the factors in their evaluation. While "\answerYes{}" is generally preferable to "\answerNo{}", it is perfectly acceptable to answer "\answerNo{}" provided a proper justification is given (e.g., "error bars are not reported because it would be too computationally expensive" or "we were unable to find the license for the dataset we used"). In general, answering "\answerNo{}" or "\answerNA{}" is not grounds for rejection. While the questions are phrased in a binary way, we acknowledge that the true answer is often more nuanced, so please just use your best judgment and write a justification to elaborate. All supporting evidence can appear either in the main paper or the supplemental material, provided in appendix. If you answer \answerYes{} to a question, in the justification please point to the section(s) where related material for the question can be found.

% IMPORTANT, please:
% \begin{itemize}
%     \item {\bf Delete this instruction block, but keep the section heading ``NeurIPS Paper Checklist"},
%     \item  {\bf Keep the checklist subsection headings, questions/answers and guidelines below.}
%     \item {\bf Do not modify the questions and only use the provided macros for your answers}.
% \end{itemize} 
 

% %%% END INSTRUCTIONS %%%


\begin{enumerate}

\item {\bf Claims}
    \item[] Question: Do the main claims made in the abstract and introduction accurately reflect the paper's contributions and scope?
    \item[] Answer: \answerYes{} % Replace by \answerYes{}, \answerNo{}, or \answerNA{}.
    \item[] Justification: The abstract and introduction summarizes motivation, objective, results and conclusion. The claims made match theoretical and experimental results. Guidelines are given in the repository on how to reproduce expected results.
    \item[] Guidelines:
    \begin{itemize}
        \item The answer NA means that the abstract and introduction do not include the claims made in the paper.
        \item The abstract and/or introduction should clearly state the claims made, including the contributions made in the paper and important assumptions and limitations. A No or NA answer to this question will not be perceived well by the reviewers. 
        \item The claims made should match theoretical and experimental results, and reflect how much the results can be expected to generalize to other settings. 
        \item It is fine to include aspirational goals as motivation as long as it is clear that these goals are not attained by the paper. 
    \end{itemize}

\item {\bf Limitations}
    \item[] Question: Does the paper discuss the limitations of the work performed by the authors?
    \item[] Answer: \answerTODO{} % Replace by \answerYes{}, \answerNo{}, or \answerNA{}.
    \item[] Justification: \justificationTODO{}
    \item[] Guidelines:
    \begin{itemize}
        \item The answer NA means that the paper has no limitation while the answer No means that the paper has limitations, but those are not discussed in the paper. 
        \item The authors are encouraged to create a separate "Limitations" section in their paper.
        \item The paper should point out any strong assumptions and how robust the results are to violations of these assumptions (e.g., independence assumptions, noiseless settings, model well-specification, asymptotic approximations only holding locally). The authors should reflect on how these assumptions might be violated in practice and what the implications would be.
        \item The authors should reflect on the scope of the claims made, e.g., if the approach was only tested on a few datasets or with a few runs. In general, empirical results often depend on implicit assumptions, which should be articulated.
        \item The authors should reflect on the factors that influence the performance of the approach. For example, a facial recognition algorithm may perform poorly when image resolution is low or images are taken in low lighting. Or a speech-to-text system might not be used reliably to provide closed captions for online lectures because it fails to handle technical jargon.
        \item The authors should discuss the computational efficiency of the proposed algorithms and how they scale with dataset size.
        \item If applicable, the authors should discuss possible limitations of their approach to address problems of privacy and fairness.
        \item While the authors might fear that complete honesty about limitations might be used by reviewers as grounds for rejection, a worse outcome might be that reviewers discover limitations that aren't acknowledged in the paper. The authors should use their best judgment and recognize that individual actions in favor of transparency play an important role in developing norms that preserve the integrity of the community. Reviewers will be specifically instructed to not penalize honesty concerning limitations.
    \end{itemize}

\item {\bf Theory assumptions and proofs}
    \item[] Question: For each theoretical result, does the paper provide the full set of assumptions and a complete (and correct) proof?
    \item[] Answer: \answerYes{} % Replace by \answerYes{}, \answerNo{}, or \answerNA{}.
    \item[] Justification: The theoretical background for the impletemented activation function is based on existing concept of geometry. Proof for the working of the implementation is provided.
    \item[] Guidelines:
    \begin{itemize}
        \item The answer NA means that the paper does not include theoretical results. 
        \item All the theorems, formulas, and proofs in the paper should be numbered and cross-referenced.
        \item All assumptions should be clearly stated or referenced in the statement of any theorems.
        \item The proofs can either appear in the main paper or the supplemental material, but if they appear in the supplemental material, the authors are encouraged to provide a short proof sketch to provide intuition. 
        \item Inversely, any informal proof provided in the core of the paper should be complemented by formal proofs provided in appendix or supplemental material.
        \item Theorems and Lemmas that the proof relies upon should be properly referenced. 
    \end{itemize}

    \item {\bf Experimental result reproducibility}
    \item[] Question: Does the paper fully disclose all the information needed to reproduce the main experimental results of the paper to the extent that it affects the main claims and/or conclusions of the paper (regardless of whether the code and data are provided or not)?
    \item[] Answer: \answerYes{} % Replace by \answerYes{}, \answerNo{}, or \answerNA{}.
    \item[] Justification: Standard datasets and exsiting state of the art models are used. The algoorithm for the Activation fucntion and settings used to reporiduce the results are provided.
    \item[] Guidelines: 
    \begin{itemize}
        \item The answer NA means that the paper does not include experiments.
        \item If the paper includes experiments, a No answer to this question will not be perceived well by the reviewers: Making the paper reproducible is important, regardless of whether the code and data are provided or not.
        \item If the contribution is a dataset and/or model, the authors should describe the steps taken to make their results reproducible or verifiable. 
        \item Depending on the contribution, reproducibility can be accomplished in various ways. For example, if the contribution is a novel architecture, describing the architecture fully might suffice, or if the contribution is a specific model and empirical evaluation, it may be necessary to either make it possible for others to replicate the model with the same dataset, or provide access to the model. In general. releasing code and data is often one good way to accomplish this, but reproducibility can also be provided via detailed instructions for how to replicate the results, access to a hosted model (e.g., in the case of a large language model), releasing of a model checkpoint, or other means that are appropriate to the research performed.
        \item While NeurIPS does not require releasing code, the conference does require all submissions to provide some reasonable avenue for reproducibility, which may depend on the nature of the contribution. For example
        \begin{enumerate}
            \item If the contribution is primarily a new algorithm, the paper should make it clear how to reproduce that algorithm.
            \item If the contribution is primarily a new model architecture, the paper should describe the architecture clearly and fully.
            \item If the contribution is a new model (e.g., a large language model), then there should either be a way to access this model for reproducing the results or a way to reproduce the model (e.g., with an open-source dataset or instructions for how to construct the dataset).
            \item We recognize that reproducibility may be tricky in some cases, in which case authors are welcome to describe the particular way they provide for reproducibility. In the case of closed-source models, it may be that access to the model is limited in some way (e.g., to registered users), but it should be possible for other researchers to have some path to reproducing or verifying the results.
        \end{enumerate}
    \end{itemize}


\item {\bf Open access to data and code}
    \item[] Question: Does the paper provide open access to the data and code, with sufficient instructions to faithfully reproduce the main experimental results, as described in supplemental material?
    \item[] Answer: \answerYes{} % Replace by \answerYes{}, \answerNo{}, or \answerNA{}.
    \item[] Justification: Code is written in Python and uses open-source free libraries. A detailed README is provided in the repo for conducting the experiments.
    \item[] Guidelines:
    \begin{itemize}
        \item The answer NA means that paper does not include experiments requiring code.
        \item Please see the NeurIPS code and data submission guidelines (\url{https://nips.cc/public/guides/CodeSubmissionPolicy}) for more details.
        \item While we encourage the release of code and data, we understand that this might not be possible, so “No” is an acceptable answer. Papers cannot be rejected simply for not including code, unless this is central to the contribution (e.g., for a new open-source benchmark).
        \item The instructions should contain the exact command and environment needed to run to reproduce the results. See the NeurIPS code and data submission guidelines (\url{https://nips.cc/public/guides/CodeSubmissionPolicy}) for more details.
        \item The authors should provide instructions on data access and preparation, including how to access the raw data, preprocessed data, intermediate data, and generated data, etc.
        \item The authors should provide scripts to reproduce all experimental results for the new proposed method and baselines. If only a subset of experiments are reproducible, they should state which ones are omitted from the script and why.
        \item At submission time, to preserve anonymity, the authors should release anonymized versions (if applicable).
        \item Providing as much information as possible in supplemental material (appended to the paper) is recommended, but including URLs to data and code is permitted.
    \end{itemize}


\item {\bf Experimental setting/details}
    \item[] Question: Does the paper specify all the training and test details (e.g., data splits, hyperparameters, how they were chosen, type of optimizer, etc.) necessary to understand the results?
    \item[] Answer: \answerYes{} % Replace by \answerYes{}, \answerNo{}, or \answerNA{}.
    \item[] Justification: Full details are provided with the code.
    \item[] Guidelines:
    \begin{itemize}
        \item The answer NA means that the paper does not include experiments.
        \item The experimental setting should be presented in the core of the paper to a level of detail that is necessary to appreciate the results and make sense of them.
        \item The full details can be provided either with the code, in appendix, or as supplemental material.
    \end{itemize}

\item {\bf Experiment statistical significance}
    \item[] Question: Does the paper report error bars suitably and correctly defined or other appropriate information about the statistical significance of the experiments?
    \item[] Answer: \answerYes{} % Replace by \answerYes{}, \answerNo{}, or \answerNA{}.
    \item[] Justification: For each activation function, experiments were repeated over seven independent noise realizations. For NLSD, performance was evaluated using the standard deviation of residuals between predicted and true values. Bars report the median across runs, and error bars indicate ±1 standard deviation across runs.
    \item[] Guidelines:
    \begin{itemize}
        \item The answer NA means that the paper does not include experiments.
        \item The authors should answer "Yes" if the results are accompanied by error bars, confidence intervals, or statistical significance tests, at least for the experiments that support the main claims of the paper.
        \item The factors of variability that the error bars are capturing should be clearly stated (for example, train/test split, initialization, random drawing of some parameter, or overall run with given experimental conditions).
        \item The method for calculating the error bars should be explained (closed form formula, call to a library function, bootstrap, etc.)
        \item The assumptions made should be given (e.g., Normally distributed errors).
        \item It should be clear whether the error bar is the standard deviation or the standard error of the mean.
        \item It is OK to report 1-sigma error bars, but one should state it. The authors should preferably report a 2-sigma error bar than state that they have a 96\% CI, if the hypothesis of Normality of errors is not verified.
        \item For asymmetric distributions, the authors should be careful not to show in tables or figures symmetric error bars that would yield results that are out of range (e.g. negative error rates).
        \item If error bars are reported in tables or plots, The authors should explain in the text how they were calculated and reference the corresponding figures or tables in the text.
    \end{itemize}

\item {\bf Experiments compute resources}
    \item[] Question: For each experiment, does the paper provide sufficient information on the computer resources (type of compute workers, memory, time of execution) needed to reproduce the experiments?
    \item[] Answer: \answerYes{} % Replace by \answerYes{}, \answerNo{}, or \answerNA{}.
    \item[] Justification: Run time per epoch, individual run and total time along with hardware spec is tabulated in Appendix.
    \item[] Guidelines:
    \begin{itemize}
        \item The answer NA means that the paper does not include experiments.
        \item The paper should indicate the type of compute workers CPU or GPU, internal cluster, or cloud provider, including relevant memory and storage.
        \item The paper should provide the amount of compute required for each of the individual experimental runs as well as estimate the total compute. 
        \item The paper should disclose whether the full research project required more compute than the experiments reported in the paper (e.g., preliminary or failed experiments that didn't make it into the paper). 
    \end{itemize}
    
\item {\bf Code of ethics}
    \item[] Question: Does the research conducted in the paper conform, in every respect, with the NeurIPS Code of Ethics \url{https://neurips.cc/public/EthicsGuidelines}?
    \item[] Answer: \answerYes{} % Replace by \answerYes{}, \answerNo{}, or \answerNA{}.
    \item[] Justification:
    \item[] Guidelines:
    \begin{itemize}
        \item The answer NA means that the authors have not reviewed the NeurIPS Code of Ethics.
        \item If the authors answer No, they should explain the special circumstances that require a deviation from the Code of Ethics.
        \item The authors should make sure to preserve anonymity (e.g., if there is a special consideration due to laws or regulations in their jurisdiction).
    \end{itemize}


\item {\bf Broader impacts}
    \item[] Question: Does the paper discuss both potential positive societal impacts and negative societal impacts of the work performed?
    \item[] Answer: \answerYes{} % Replace by \answerYes{}, \answerNo{}, or \answerNA{}.
    \item[] Justification: The proposed activation function may have positive impacts by improving model efficiency, stability, or performance across machine learning applications. As a low level mathematical component, it does not introduce direct societal risks; any downstream negative impacts would stem from how machine learning systems are applied rather than from the activation function itself.
    \item[] Guidelines:
    \begin{itemize}
        \item The answer NA means that there is no societal impact of the work performed.
        \item If the authors answer NA or No, they should explain why their work has no societal impact or why the paper does not address societal impact.
        \item Examples of negative societal impacts include potential malicious or unintended uses (e.g., disinformation, generating fake profiles, surveillance), fairness considerations (e.g., deployment of technologies that could make decisions that unfairly impact specific groups), privacy considerations, and security considerations.
        \item The conference expects that many papers will be foundational research and not tied to particular applications, let alone deployments. However, if there is a direct path to any negative applications, the authors should point it out. For example, it is legitimate to point out that an improvement in the quality of generative models could be used to generate deepfakes for disinformation. On the other hand, it is not needed to point out that a generic algorithm for optimizing neural networks could enable people to train models that generate Deepfakes faster.
        \item The authors should consider possible harms that could arise when the technology is being used as intended and functioning correctly, harms that could arise when the technology is being used as intended but gives incorrect results, and harms following from (intentional or unintentional) misuse of the technology.
        \item If there are negative societal impacts, the authors could also discuss possible mitigation strategies (e.g., gated release of models, providing defenses in addition to attacks, mechanisms for monitoring misuse, mechanisms to monitor how a system learns from feedback over time, improving the efficiency and accessibility of ML).
    \end{itemize}
    
\item {\bf Safeguards}
    \item[] Question: Does the paper describe safeguards that have been put in place for responsible release of data or models that have a high risk for misuse (e.g., pretrained language models, image generators, or scraped datasets)?
    \item[] Answer: \answerNA{} % Replace by \answerYes{}, \answerNo{}, or \answerNA{}.
    \item[] Justification: This work proposes a mathematical activation function and does not introduce datasets, pretrained models, or artifacts that pose a risk of misuse. As such, no special safeguards for responsible release are required.
    \item[] Guidelines:
    \begin{itemize}
        \item The answer NA means that the paper poses no such risks.
        \item Released models that have a high risk for misuse or dual-use should be released with necessary safeguards to allow for controlled use of the model, for example by requiring that users adhere to usage guidelines or restrictions to access the model or implementing safety filters. 
        \item Datasets that have been scraped from the Internet could pose safety risks. The authors should describe how they avoided releasing unsafe images.
        \item We recognize that providing effective safeguards is challenging, and many papers do not require this, but we encourage authors to take this into account and make a best faith effort.
    \end{itemize}

\item {\bf Licenses for existing assets}
    \item[] Question: Are the creators or original owners of assets (e.g., code, data, models), used in the paper, properly credited and are the license and terms of use explicitly mentioned and properly respected?
    \item[] Answer: \answerYes{} % Replace by \answerYes{}, \answerNo{}, or \answerNA{}.
    \item[] Justification: models used are cited to their original papers. Standard datasets are used without re-packaging.
    \item[] Guidelines:
    \begin{itemize}
        \item The answer NA means that the paper does not use existing assets.
        \item The authors should cite the original paper that produced the code package or dataset.
        \item The authors should state which version of the asset is used and, if possible, include a URL.
        \item The name of the license (e.g., CC-BY 4.0) should be included for each asset.
        \item For scraped data from a particular source (e.g., website), the copyright and terms of service of that source should be provided.
        \item If assets are released, the license, copyright information, and terms of use in the package should be provided. For popular datasets, \url{paperswithcode.com/datasets} has curated licenses for some datasets. Their licensing guide can help determine the license of a dataset.
        \item For existing datasets that are re-packaged, both the original license and the license of the derived asset (if it has changed) should be provided.
        \item If this information is not available online, the authors are encouraged to reach out to the asset's creators.
    \end{itemize}

\item {\bf New assets}
    \item[] Question: Are new assets introduced in the paper well documented and is the documentation provided alongside the assets?
    \item[] Answer: \answerNA{} % Replace by \answerYes{}, \answerNo{}, or \answerNA{}.
    \item[] Justification:
    \item[] Guidelines:
    \begin{itemize}
        \item The answer NA means that the paper does not release new assets.
        \item Researchers should communicate the details of the dataset/code/model as part of their submissions via structured templates. This includes details about training, license, limitations, etc. 
        \item The paper should discuss whether and how consent was obtained from people whose asset is used.
        \item At submission time, remember to anonymize your assets (if applicable). You can either create an anonymized URL or include an anonymized zip file.
    \end{itemize}

\item {\bf Crowdsourcing and research with human subjects}
    \item[] Question: For crowdsourcing experiments and research with human subjects, does the paper include the full text of instructions given to participants and screenshots, if applicable, as well as details about compensation (if any)? 
    \item[] Answer: \answerNA{} % Replace by \answerYes{}, \answerNo{}, or \answerNA{}.
    \item[] Justification:
    \item[] Guidelines:
    \begin{itemize}
        \item The answer NA means that the paper does not involve crowdsourcing nor research with human subjects.
        \item Including this information in the supplemental material is fine, but if the main contribution of the paper involves human subjects, then as much detail as possible should be included in the main paper. 
        \item According to the NeurIPS Code of Ethics, workers involved in data collection, curation, or other labor should be paid at least the minimum wage in the country of the data collector. 
    \end{itemize}

\item {\bf Institutional review board (IRB) approvals or equivalent for research with human subjects}
    \item[] Question: Does the paper describe potential risks incurred by study participants, whether such risks were disclosed to the subjects, and whether Institutional Review Board (IRB) approvals (or an equivalent approval/review based on the requirements of your country or institution) were obtained?
    \item[] Answer: \answerNA{} % Replace by \answerYes{}, \answerNo{}, or \answerNA{}.
    \item[] Justification:
    \item[] Guidelines:
    \begin{itemize}
        \item The answer NA means that the paper does not involve crowdsourcing nor research with human subjects.
        \item Depending on the country in which research is conducted, IRB approval (or equivalent) may be required for any human subjects research. If you obtained IRB approval, you should clearly state this in the paper. 
        \item We recognize that the procedures for this may vary significantly between institutions and locations, and we expect authors to adhere to the NeurIPS Code of Ethics and the guidelines for their institution. 
        \item For initial submissions, do not include any information that would break anonymity (if applicable), such as the institution conducting the review.
    \end{itemize}

\item {\bf Declaration of LLM usage}
    \item[] Question: Does the paper describe the usage of LLMs if it is an important, original, or non-standard component of the core methods in this research? Note that if the LLM is used only for writing, editing, or formatting purposes and does not impact the core methodology, scientific rigorousness, or originality of the research, declaration is not required.
    %this research? 
    \item[] Answer: \answerNA{} % Replace by \answerYes{}, \answerNo{}, or \answerNA{}.
    \item[] Justification:
    \item[] Guidelines:
    \begin{itemize}
        \item The answer NA means that the core method development in this research does not involve LLMs as any important, original, or non-standard components.
        \item Please refer to our LLM policy (\url{https://neurips.cc/Conferences/2025/LLM}) for what should or should not be described.
    \end{itemize}

\end{enumerate}


\end{document}